\subsection{window-\texorpdfstring{$6$}{6} 的 anomaly 相位预算刚性与连续承载维数}\label{subsec:conclusion-window6-phase-budget-rigidity}

沿用定理 \ref{thm:xi-foldbin-parity-charge-split-exact-minrank}、
命题 \ref{prop:conclusion-window6-quadratic-collision-mass-anomaly-discriminant-closure}
与猜想 \ref{conj:conclusion-window6-anomaly-collision-discriminant-factorization} 的记号。前一小节已经把
\[
\deg\Delta_6=21+53
\]
压缩成 anomaly 账本与二次碰撞质量的数值闭合，但这一分裂仍可进一步收紧为一个类型刚性结论：其中 \(53\) 只给出标量化的碰撞质量，而 \(21\) 则不是任意可再压缩的计数，它恰是一个基本二群在连续相位几何中所需的最小忠实承载维数。因此，window-\(6\) 的 \(21+53\) 不是同类资源的加法，而是连续相位预算与碰撞统计预算的严格类型分裂。

\begin{theorem}[基本二群异常签名的最小连续相位承载维数]\label{thm:conclusion-elementary-2group-minimal-torus-dimension}
设
\[
A\cong (\ZZ/2\ZZ)^\nu
\]
为有限基本二群。则
\[
\min\bigl\{r:\exists\ \text{单射群同态 }A\hookrightarrow \mathbb T^r\bigr\}=\nu.
\]
亦即：一个 \(\nu\)-维二元异常签名，恰需要 \(\nu\) 个连续相位坐标才能被忠实承载；少一个维数都不可能。
\end{theorem}

\begin{proof}
任意群同态 \(A\to \mathbb T^r\) 的像都由 \(2\)-扭元组成，因为对任意 \(a\in A\) 都有 \(2a=0\)。故其像必包含于
\[
\mathbb T^r[2]:=\{z\in \mathbb T^r:2z=0\}\cong (\ZZ/2\ZZ)^r.
\]
若该同态单射，则必须存在单射
\[
(\ZZ/2\ZZ)^\nu\hookrightarrow (\ZZ/2\ZZ)^r,
\]
从而必有 \(\nu\le r\)。

反之，当 \(r=\nu\) 时，取坐标符号嵌入
\[
(\epsilon_1,\dots,\epsilon_\nu)\longmapsto
\bigl((-1)^{\epsilon_1},\dots,(-1)^{\epsilon_\nu}\bigr)\in \mathbb T^\nu
\]
即可得到一个单射群同态。故最小维数正是 \(\nu\)。证毕。
\end{proof}

\begin{corollary}[window-\texorpdfstring{$6$}{6} 的 anomaly 相位预算刚性与 collision 预算的类型不可交换性]\label{cor:conclusion-window6-anomaly-phase-budget-rigidity}
在 window-\(6\) 中，由命题 \ref{prop:conclusion-window6-quadratic-collision-mass-anomaly-discriminant-closure}，
\[
\deg\Delta_6
=
\dim_{\FF_2}\bigl(G_6^{\mathrm{bin}}\bigr)^{\mathrm{ab}}
+\mathfrak c_6
=
21+53
=
74.
\]
因此，window-\(6\) 的判别复杂度严格分裂为两部分：一部分是 \(21\) 维 anomaly 账本，另一部分是 \(53\) 单位的二次碰撞质量。并且前一部分若要以连续相位几何忠实承载，则至少需要
\[
\mathbb T^{21}
\]
这一维数；任何 \(r<21\) 的 torus 都不可能忠实实现该异常部分。于是
\[
74=21+53
\]
不是纯计数分解，而是类型分裂：\(21\) 是不可压缩的连续相位预算，\(53\) 是碰撞统计预算，两者不能彼此兑换。
\end{corollary}

\begin{proof}
由定理 \ref{thm:xi-foldbin-parity-charge-split-exact-minrank}，
\[
\bigl(G_6^{\mathrm{bin}}\bigr)^{\mathrm{ab}}\cong (\ZZ/2\ZZ)^{21}.
\]
对该基本二群应用定理 \ref{thm:conclusion-elementary-2group-minimal-torus-dimension}，即可得到：任何连续 torus 承载若要对 anomaly 部分保持忠实，维数必至少为 \(21\)。

另一方面，命题 \ref{prop:conclusion-window6-quadratic-collision-mass-anomaly-discriminant-closure} 已给出
\[
\mathfrak c_6=53,
\qquad
\deg\Delta_6=21+53.
\]
若试图把 \(21\) 维 anomaly 预算并入 \(53\) 单位碰撞质量并同时降低所需 torus 维数，则将得到一个到 \(\mathbb T^r\)（\(r<21\)）的忠实表示，这与定理 \ref{thm:conclusion-elementary-2group-minimal-torus-dimension} 矛盾。故 anomaly 相位预算与 collision 统计预算在类型上不可互换。证毕。
\end{proof}

\paragraph{统一读法}
至此，window-\(6\) 的判别复杂度可被压成三层互锁对象：\(\ZZ/(21\ZZ)\) 的可见 CRT 相空间负责算术可见层，\(A_1^8A_2^4A_3^9\) 的隐藏反射包负责几何与判别层，而
\[
\bigl(G_6^{\mathrm{bin}}\bigr)^{\mathrm{ab}}\cong (\ZZ/2\ZZ)^{21}
\]
则把其中不可约的异常部分钉死为一个最小连续承载维数恰等于 \(21\) 的二元相位账本。故 \(21+53\) 不是把同一种复杂度拆成两块，而是把连续相位缺陷与碰撞统计缺陷区分为两种不可互换的资源。

\endinput
